\clearpage
\thispagestyle{empty}
\pdfbookmark[0]{Modern edition: authorship and sources}{modern-edition-notice}
\begingroup
\pagestyle{empty}
\footnotesize\raggedright
\enlargethispage{2\baselineskip}
\setlength{\parindent}{0pt}\setlength{\parskip}{.5\baselineskip}
{\large Modern edition: authorship and sources\par}
Historical author: Carl Friedrich Gauss. The historical editors, correspondents, contributors and issuing institutions retain the attributions printed in the source and reproduced in the historical text. Their writing is not attributed to an AI system.

Source edition: Carl Friedrich Gauss, \emph{Werke}, Band VII, 1906, containing \emph{Theoria motus corporum coelestium in sectionibus conicis solem ambientium}, Hamburg, 1809. Digital witness: \href{https://gdz.sub.uni-goettingen.de/id/PPN236008730}{Göttinger Digitalisierungszentrum, PPN236008730}.

This original-language partial edition contains the historical titles, complete Praefatio and Contenta, Articles 1--\GaussCurrentLastArticle, the Tabulae part-title and complete Tables I--\GaussCurrentLastTable, Figures 1--\GaussCurrentFigureCount, \GaussCurrentBemerkungenScope, and the Zusätze part-title, Nachlass heading and complete items 1--\GaussCurrentLastZusatz, through [18.], \GaussCurrentZusatzRemarksScope, including Oriani’s corrections, the historical note and reply by Gauss, the ninth and tenth items’ geometric proofs and editable orbital and spherical diagrams, the eleventh item’s historical correction notice, the twelfth item’s determination of probable errors in least-squares results, the thirteenth item’s relations between initial state and orbital elements, the fourteenth item’s correction of planetary elements from four oppositions, the fifteenth item’s differential formulas with a Pallas example, and the sixteenth item’s discussion of celestial zodiacs with Goldschmidt’s Iris table. The part-title Zur parabolischen Bewegung, Nachlass und Briefwechsel heading and complete notes 1--\GaussCurrentLastParabolicNote{} on elapsed time from two radius vectors and their chord, the inverse calculation of the chord, determination of the orbital elements, the triangle-to-sector relation with its logarithm table, and Gauss’s letters to Olbers of 7 and 13 January 1815 with a worked calculation and the subsequent reduction of the cometary-orbit problem, followed by the general theory of cometary-orbit calculation and its full example for the second comet of 1813, and the letters to Olbers of 29 May and to Bessel of 24 June 1815 concerning the proposed supplement and parabolic-motion table, are also included. The Nachlass und Veröffentlichung heading, the second part on the true anomaly in parabolic motion, and explanatory Articles 1--\GaussCurrentLastTabulaNovaArticle{} of Tabula nova motus parabolici are included, completing its explanation, together with \GaussCurrentTabulaNovaTableScope. \GaussCurrentPostTabulaNovaScope{} are also included. \GaussCurrentCeresScope{} are also included. The remaining contents of Volume VII are outside this edition. Historical wording, mathematical expressions and notes are retained. Modern fonts and reflow represent content rather than historical type identity. Qualified readings are described in the accompanying README.

AI-transcribed text, AI-produced source comparisons and typesetting, and AI-created editable mathematical drawings: OpenAI GPT-6 Astra (Ultra). This credit concerns the modern encoding and does not replace historical authorship. No human or independent review is claimed.

In Tabula nova, Pars II, printed page 359, the final digit of argument 5,354 and the zero in 20,228 at argument 5,294 are qualified readings from incomplete contours. The point above the final 6 of the value at 5,012 records an observed mark without assigning mathematical meaning; its printing intent is uncertain. On printed page 360, the filled superscript mark after 67 in the second block’s first row and the final digit of argument 5,620 are read as zero with support from corresponding forms and the comparison witness; their principal interiors are not fully visible. In Pars III, printed page 367, the seconds digit at argument 32200 is read as 8 from a surviving lower loop and crossing stroke; its upper loop is incomplete and the reading remains qualified. In the Encke letter, printed page 368, a lowered colon represents two vertically separated contours beside the first denominator A. Their punctuation identity and printing intent are uncertain; damaged type or printing/imaging irregularity remain alternatives. No mathematical meaning is assigned to the mark. The comma following the logarithmic interpolation fraction on printed page 369 is represented at its observed fraction-bar height. In the letter of 1 April 1843, the first calculation on printed page 372 reads the dense exponent over q as 3/2 with qualified confidence; the clearer earlier definition supports this reading, while a faint comparison witness does not resolve every stroke. A lowered period after Gleichung on printed page 373 represents a detached point at descender height; its punctuation identity and printing intent remain uncertain. In the editorial Bemerkungen on printed page 374, the comma following the sector-to-triangle fraction is represented at its observed fraction-bar height; its printing intent remains uncertain.
 In Ceres section [2], printed page 378, the first denominator exponent in the definition of R is read as 3 with support from the comparison witness; its dense principal contour also admits a reading of 2 or filled/damaged type. The final prime and closing parenthesis in the second root are qualified readings from bent and fragmentary contours. The raised formula punctuation on printed pages 378–381 represents observed positions without added mathematical meaning; printing intent remains uncertain.
 In Ceres section [3], printed page 382, the two raised commas retain their observed positions. A plus with a detached lower stroke is represented by a normal editable plus without a claim of continuous ink; the cause of the interruption is uncertain. The printed unprimed M in the Delta-w definition is retained.
 In Ceres section [4], the point above the equals sign in the first geometric-mean relation and the low point after P in the C relation retain their observed positions without added mathematical meaning. Their printing intent is uncertain. The raised comma after the limiting P expression and the printed root with following fractional exponent are retained.
 In Ceres section [5], the final M and closing parenthesis of the first integrated product are supplied by the Internet Archive witness, native page 397; the corresponding GDZ expression ends visibly at the preceding 2. The witness parenthesis survives as a lower curved stroke, so its reading is qualified. The printed coefficients 35/24 and 11/8 and the expression 3entM are retained without mathematical emendation.
 In Ceres section [6], the two terminal M letters in the second velocity row are supplied by the Internet Archive witness, native page 399; both endings are lost at the right edge of the GDZ image. The adjacent closing parentheses and final brace are qualified reconstructed readings: surviving fragments, the opening delimiters and the recovered expression support them, but the complete glyph contours are not resolved. Ordinary and Fraktur B, the unprimed e in the fourth velocity result, all 98 leader points and six underlines follow the principal without mathematical emendation.
 In Ceres section [7], the first f coefficient is 1/24, and the integration constants retain the printed primes. One horizontal mark above a primed tilded omega is represented by an additional bar without assigning it mathematical meaning. The a-like Greek alpha reading has moderate character confidence: its low right foot differs from the curled Latin a, while the recovered reading and its pairing with beta provide supporting context. Several omega-primes contact closing-parenthesis contours; modern separate glyphs do not assert separate historical ink. Raised terminal punctuation and twelve individual leader points retain their observed positions.
 In Ceres section [8], a compact detached point below and to the right of the first sine identity is represented by a lowered period; its punctuation identity and printing intent remain uncertain. The terminal point of the complete periodic-latitude formula is represented at fraction-bar height without a multiplication meaning. A small irregular rising contour below the minus in XI is not confidently a separate sign. Incidental ink, damaged type, printing or imaging irregularity remain possible.
 In the Ceres perturbation-table letters, short horizontal strokes above the left of N in Nov. and above the right of n in angezeigten are represented by editable short rules. Raised points between die and VII. and after und before C,D retain their observed positions. These readings have high contour and position confidence but uncertain punctuation identity or printing intent; no additional linguistic or mathematical meaning is assigned. Irregular contours near ist, als, und in the formula array, gehen and Methode are not confidently separate signs. Incidental ink, damaged type, printing or imaging irregularity remain possible. The Ceres, Jupiter, Pallas, opposition and ascending-node symbols retain their distinct identities; modern fonts do not reproduce historical type.
 In the letter of 10 May 1805, the final astronomical symbols are Pallas and Juno, read from their pointed-loop and radiating-star heads with lower crosses. The opening parenthesis before denn survives in separated upper and lower curved strokes; its editable form records the recognizable glyph without asserting continuous ink. A tiny partly hollow angular contour above the left of w in werde is not confidently a separate linguistic sign. Contour and position confidence is high; the cause of the parenthesis interruption and the detached mark's printing intent remain uncertain. Incidental ink, damaged type, printing or imaging irregularity remain possible.
 In section [2] of the second Ceres method, the first reciprocal relation on printed page 403 has the fractional exponent encoded as 8/2 from an 8-like contour, with moderate numeral confidence. The Internet Archive comparison, leaf 413, also appears 8-like but does not resolve every stroke. A filled or damaged 3 is an alternative supported by the preceding mathematical expression and the retained transcription; that contextual inference is not substituted for the qualified visual reading. The two parentheses around S. 381 have interrupted contours. The derivative comma and reciprocal-relation point retain their observed fraction-bar heights. Detached marks below entwickeln and below/right of the final differential are not confidently separate signs. Their printing intent and the cause of interrupted contours remain uncertain; incidental ink, damaged type, printing or imaging irregularity remain possible.
 In section [3] of the second Ceres method, the comma following the constant fraction retains its observed fraction-bar height without added mathematical meaning. A compact oblique mark before vgl in the historical note, a shallow two-winged contour above the first odd-case minus and a compact point well to the right of Formeln are not confidently separate signs. Their contours and positions are clear, but their printing intent remains uncertain; incidental ink, damaged type, printing or imaging irregularity remain possible.
 In section [4] of the second Ceres method, the beta-fraction semicolon retains its observed fraction-bar position. A compact angular contour after the logr-prime table header and a tiny branched contour below the final printed-page405 formula are not confidently separate signs; incidental ink, damaged type, printing or imaging irregularity remain possible. The last first-row expression on the right of the wide printed-page406 table is transcribed through +M without an added closing parenthesis: that character is not confidently resolved at the principal edge, and the faint corresponding edge of the retained Internet Archive witness does not settle it. This is a qualified boundary reading, not a claim of a historical printing error or absence of ink. The wide table uses a landscape page; its shared row associations and both complete braced headers remain intact.
 In section [5] of the second Ceres method, the first opening parenthesis has interrupted curved strokes; the editable glyph represents the recognizable reading without asserting continuous ink. A small flattened angular contour well to the right of the first printed-page407 numerical row is not confidently separate punctuation. Incidental ink, damaged type, printing or imaging irregularity remain possible.
 In the letter of 2 July 1805, interrupted curved contours in formula A and the four-row group in B are represented by recognizable closing parentheses and a grouping brace without asserting continuous ink. Detached angular contours above-left of the i-dot in the +3,59 sine term and below-left of the node in the -4,56 term are not confidently separate signs. Incidental ink, damaged type, printing or imaging irregularity remain possible. All eighteen underlines, eight grouping braces, sixty-four individual leader points and the printed numerical distinctions are retained without mathematical emendation.
 In Brendel’s closing remarks on Ceres, a detached angular contour above the o/r gap in the second occurrence of Correspondenz is not confidently a separate linguistic sign. Incidental ink, damaged type, printing or imaging irregularity remain possible. Brendel’s account of his reconstruction from Gauss’s numerical work is retained in full.

\par\medskip
{\large Modern edition: comparison witnesses\par}
Alternative witness for Articles 98, 124 and 179, Table II and the sixth, seventh, eighth, eleventh, thirteenth and fifteenth Zusätze: \href{https://archive.org/details/werkehrsgvonderg07gausuoft}{Internet Archive, \emph{Werke}, volume 7, item werkehrsgvonderg07gausuoft}, Gerstein--University of Toronto digitization; PDF \texttt{werkehrsgvonderg07gausuoft.pdf}, leaves 140--141, printed pages 130--131, for Article 98; leaf 175, printed page 165, for Article 124. For Article 98 this witness supports the qualified exponent reading and comparison of small detached marks. For Article 124 it supplies a comparison for the qualified unprimed reading of \(X\) beside a small detached contour in the principal witness. For Article 179, leaf 255, printed page 245, supplies a faint comparison for the partially surviving exponent parentheses; it does not independently resolve their complete contours. For Table II, leaves 280 and 285, printed pages 270 and 275, support qualified readings in the historical note and two numerical entries. For the sixth Zusatz, leaf 309, printed page 299, supports the qualified opening-fraction reading and the recurrence coefficients. For the seventh Zusatz, leaves 310--311, printed pages 300--301, supply comparisons for the dense nested coefficients and final exponent; their internal contours remain partly unresolved. For the eighth Zusatz, leaf 312, printed page 302, supports the reading of the damaged closing parenthesis and double-prime in the second row of the twelve-term expression. For the eleventh Zusatz, leaf 317, printed page 307, supplies a comparison for the fragmented bracket-shaped contour after und; it does not corroborate the qualified principal reading. For the thirteenth Zusatz, leaf 319, printed page 309, supports the qualified exponent reading 3/2 in the final integral; dense principal contours remain partly unresolved. Leaves 321–322, printed pages 311–312, also support the qualified variable-glyph readings in the fifteenth Zusatz. For the first parabolic-motion note, leaf 336, printed page 326, supplies a faint comparison for the qualified exponent reading 2 in formula [II]; it does not independently resolve that digit. For the fourth parabolic-motion note, leaf 339, printed page 329, supplies a faint comparison for the qualified first exponent reading 3/2; it does not independently resolve the dense numerator. For the sixth parabolic-motion letter, leaves 345–346, printed pages 335–336, supply comparisons for the initial fraction, root parentheses, later cosine cube, historical-note fractions and the incomplete denominator letter; the latter remains unresolved in this faint witness. For the seventh parabolic-motion note, leaves 352, 355–356 and 359, printed pages 342, 345–346 and 349, supply comparisons for the dense cubic and fractional exponents, the first series coefficient and two damaged negative indices. For the third explanatory article of Tabula nova motus parabolici, leaf 363, printed page 353, supplies comparisons for the dense fourth powers and fractional range exponents; the latter remain partly unresolved in both witnesses. For the fifth explanatory article, leaves 364–365, printed pages 354–355, supply comparisons for six dense exponent or fractional-denominator readings of 3. These remain qualified readings supported by glyph comparison and mathematical context; the faint alternative does not resolve every contour. For the second numerical part of Tabula nova, leaf 369, printed page 359, supplies comparisons for the fragmentary final digit in argument 5,354 and the fragmentary zero in the value at 5,294; neither is independently resolved in this faint witness. Leaf 370, printed page 360, supplies comparisons for two dense zero marks, the first 9 in the value at 5,524, and the fractional digits at 5,704; the faint images do not independently resolve every contour. For the third numerical part, leaf 377, printed page 367, supplies a faint comparison for the qualified 8 at argument 32200; it does not independently resolve the incomplete upper loop. For the Encke letter, leaf 378, printed page 368, supplies a faint comparison for the two-component mark beside the first denominator A; it does not resolve its printing intent. For the letter of 1 April 1843, leaf 382, printed page 372, supplies the faint exponent comparison. For Ceres section [2], leaf 388, printed page 378, supports the qualified cubic denominator reading and supplies a comparison for the bent last prime and fragmented closing parenthesis of the second root. The faint comparison images do not independently resolve every contour. Physical independence of the typesetting or image processing is not asserted.

These editions are freely available for study and teaching.

\par\clearpage\endgroup
